% titel: Hohlzylinder-Pendel (physisches Pendel)
% skript: Kap. Schwingungen, physisches Pendel, Steinerscher Satz

Ein homogener Hohlzylinder (Innenradius $r$, Aussenradius $R$, Masse $m$)
hängt auf einem (dickenlosen) Stab, der um $r$ parallel zur Symmetrieachse
versetzt ist. Dadurch entsteht ein physisches Pendel um die Drehachse $D$.
Gefragt sind die Bewegungsgleichung, das Trägheitsmoment $I_D$ und die
Eigenfrequenz.

\begin{werkzeugkasten}[Werkzeuge]
\textbf{Physisches Pendel:} Für einen um eine Achse $D$ drehbaren Körper mit
Trägheitsmoment $I_D$, dessen Schwerpunkt bei Winkel $\varphi$ unterhalb/über
der Achse hängt (Abstand $s$ von $D$ zur Schwerpunktslinie), liefert der
Drehimpulssatz
\begin{equation}
  I_D\,\ddot\varphi = -m\,g\,s\,\sin\varphi
  \;\Rightarrow\;
  \ddot\varphi + \frac{m g s}{I_D}\sin\varphi = 0.
\end{equation}

\textbf{Kleinschwingung:} $\sin\varphi\approx\varphi$:
\begin{equation}
  \ddot\varphi + \omega_0^2\,\varphi=0,
  \qquad \omega_0 = \sqrt{\frac{m g s}{I_D}}.
\end{equation}

\textbf{Steinerscher Satz} (Trägheitsmoment um eine zu einer
Symmetrieachse parallel verschobene Achse im Abstand $d$):
\begin{equation}
  I_D = I_S + m\,d^2.
\end{equation}

\textbf{Hohlzylinder} (homogen, Massen $m$, Innenradius $r$, Aussenradius $R$):
\begin{equation}
  I_S = \frac12 m\,(R^2+r^2)
\end{equation}
um die Symmetrieachse durch den Schwerpunkt.

\textbf{Konstante:} $g=\qty{9.81}{m/s^2}$.
\end{werkzeugkasten}

\begin{werkzeugkasten}[Vorgehen in allgemeiner Form]
(a) Drehimpulssatz um $D$ aufstellen (Schwerpunktabstand $s=r$); (b)
$I_D=I_S+m r^2$ mit $I_S=\frac12 m(R^2+r^2)$; (c) $\omega_0$; (d)
$\omega_0^2(r)$ für konstant gehaltenes $R,m$ maximieren.
\end{werkzeugkasten}

\begin{achtung}[Typische Fallen]
\begin{itemize}
  \item $s$ ist der Abstand \emph{Drehachse $D$--Schwerpunkt}; hier $s=r$.
  \item Der Satz von Steiner versetzt die Achse \emph{parallel}, nicht die
       Flächenmomente in einem Integral mischen.
  \item Maximierung: $\omega_0^2\propto r/(R^2+3r^2)$ bei $R=$\,konst. Nach
       $r$ ableiten.
\end{itemize}
\end{achtung}
